\subsubsection{R Utils}
\begin{enumerate}
    \item \textbf{R BST Interface}
          
          \begin{enumerate}
            \item \textbf{Build BST Index}
                  
                  R interface for building BST index. Uses R's native 1-based indexing.
                  
                  \textbf{Input:}
                  \begin{itemize}
                    \item \texttt{x numeric}: Vector of real values
                  \end{itemize}
                  \textbf{Output:}
                  \begin{itemize}
                    \item \texttt{integer}: Permutation indices (1-based R indices)
                  \end{itemize}
                  \begin{lstlisting}[language=R]
ix <- build_bst_index(x)
                  \end{lstlisting}
                  
            \item \textbf{BST Range Query}
                  
                  Performs range queries on BST index. Returns a list with components.
                  
                  \textbf{Input:}
                  \begin{itemize}
                    \item \texttt{x numeric}: Original values vector
                          
                    \item \texttt{ix integer}: BST index vector from build\_bst\_index (1-based)
                          
                    \item \texttt{lo numeric}: Lower bound of range (inclusive)
                          
                    \item \texttt{hi numeric}: Upper bound of range (inclusive)
                  \end{itemize}
                  \textbf{Output:}
                  \begin{itemize}
                    \item \texttt{list}: Contains indices and count components
                  \end{itemize}
                  \begin{lstlisting}[language=R]
result <- bst_range_query(x, ix, lo, hi)
indices <- result$indices  # 1-based indices
count <- result$count
                  \end{lstlisting}
                  
            \item \textbf{Get Sorted Value}
                  
                  Retrieves value from sorted position using R 1-based indexing.
                  
                  \textbf{Input:}
                  \begin{itemize}
                    \item \texttt{x numeric}: Original values vector
                          
                    \item \texttt{ix integer}: BST index vector (1-based)
                          
                    \item \texttt{position integer}: Position in sorted order (1-based)
                  \end{itemize}
                  \textbf{Output:}
                  \begin{itemize}
                    \item \texttt{numeric}: Value at specified position
                  \end{itemize}
                  \begin{lstlisting}[language=R]
value <- get_sorted_value(x, ix, position)
                  \end{lstlisting}
          \end{enumerate}
          
    \item \textbf{R k-d Tree Interface}
          
          \begin{enumerate}
            \item \textbf{Build k-d Tree Index}
                  
                  R interface for k-d tree construction. Input matrix dimensions are
                  rows=dimensions, columns=points.
                  
                  \textbf{Input:}
                  \begin{itemize}
                    \item \texttt{X matrix}: Matrix of points (nrow = dimensions, ncol =
                          n\_points)
                          
                    \item \texttt{dim\_order integer}: Dimension ordering (1-based, optional)
                  \end{itemize}
                  \textbf{Output:}
                  \begin{itemize}
                    \item \texttt{integer}: Permutation indices (1-based R indices)
                  \end{itemize}
                  \begin{lstlisting}[language=R]
kd_ix <- build_kd_index(X, dim_order)
                  \end{lstlisting}
                  
            \item \textbf{Build Spherical k-d Tree}
                  
                  R interface for spherical k-d tree construction with separate implementation.
                  
                  \textbf{Input:}
                  \begin{itemize}
                    \item \texttt{V matrix}: Matrix of vectors (nrow = dimensions, ncol =
                          n\_vectors)
                          
                    \item \texttt{dim\_order integer}: Dimension ordering (1-based, optional)
                  \end{itemize}
                  \textbf{Output:}
                  \begin{itemize}
                    \item \texttt{integer}: Permutation indices (1-based R indices)
                  \end{itemize}
                  \begin{lstlisting}[language=R]
sphere_ix <- build_spherical_kd(V, dim_order)
                  \end{lstlisting}
                  
            \item \textbf{Get Point from k-d Index}
                  
                  Retrieves point from k-d tree using R matrix indexing.
                  
                  \textbf{Input:}
                  \begin{itemize}
                    \item \texttt{X matrix}: Original points matrix
                          
                    \item \texttt{kd\_ix integer}: k-d tree index vector (1-based)
                          
                    \item \texttt{position integer}: Position in k-d tree order (1-based)
                  \end{itemize}
                  \textbf{Output:}
                  \begin{itemize}
                    \item \texttt{numeric}: Point values as vector
                  \end{itemize}
                  \begin{lstlisting}[language=R]
point <- get_kd_point(X, kd_ix, position)
                  \end{lstlisting}
          \end{enumerate}
\end{enumerate}